% LEGACY DOCUMENT (historical data paper wave 24)
% Current status: See RESEARCH_STATUS.md and TECH_TREE.md for canonical assessment.
% Claim status: high_risk_or_falsified (wave24-phi-ablation-falsified, Layer L5)

%%%%%%%%%%%%%%%%%%%%%%%%%%%%%%%%%%%%%%%%%%%%%%%%%%%%%%%%%%%%%%%%%%%%%%%%%%%%%%%
%  IGLA RACE φ-Ablation Falsification Report — Wave 24
%
%  Data paper — documenting a controlled ablation test on the IGLA RACE
%  benchmark.  Not a theory paper; makes no derivation claims.
%  Version: Wave 24
%%%%%%%%%%%%%%%%%%%%%%%%%%%%%%%%%%%%%%%%%%%%%%%%%%%%%%%%%%%%%%%%%%%%%%%%%%%%%%%

\documentclass[11pt,a4paper]{article}

% ─── Packages ─────────────────────────────────────────────────────────────────
\usepackage[utf8]{inputenc}
\usepackage[T1]{fontenc}
\usepackage{amsmath,amssymb,amsthm}
\usepackage{booktabs}
\usepackage{longtable}
\usepackage{graphicx}
\usepackage{hyperref}
\usepackage{url}
\usepackage[numbers,sort&compress]{natbib}
\usepackage{caption}
\usepackage{xcolor}

% ─── Hyperref setup ───────────────────────────────────────────────────────────
\hypersetup{
    colorlinks=true,
    linkcolor=blue,
    citecolor=blue,
    urlcolor=blue,
    pdftitle={IGLA RACE φ-Ablation Falsification Report — Wave 24},
    pdfsubject={Data paper, machine learning phenomenology},
    pdfkeywords={falsification, ablation, golden ratio, hyperparameters, language model, BPB}
}

% ─── Custom commands ──────────────────────────────────────────────────────────
\newcommand{\ph}{\varphi}

% ─── Begin Document ───────────────────────────────────────────────────────────
\begin{document}

% ═══════════════════════════════════════════════════════════════════════════════
%  TITLE AND AUTHORS
% ═══════════════════════════════════════════════════════════════════════════════
\title{IGLA RACE $\varphi$-Ablation Falsification Report — Wave~24}
\author{[Author List TBD]\\
  \texttt{trinity-s3ai@proton.me}}
\date{June 2026}
\maketitle

% ═══════════════════════════════════════════════════════════════════════════════
%  ABSTRACT
% ═══════════════════════════════════════════════════════════════════════════════
\begin{abstract}
We present a controlled ablation test on the IGLA RACE benchmark that
falsifies the hypothesis that $\varphi$-anchored hyperparameters improve
language-model Bits Per Byte (BPB).  The treatment arm uses
learning-rate $\varphi^{-3}\approx 0.236$, weight-decay
$\varphi^{-3}\approx 0.236$, and $\beta_{1}=1/\varphi\approx 0.618$
(AdamW, bf16, 81\,K steps, single seed 1597, 2-layer Transformer,
hidden=128, character-level modeling).  The conventional control arm
uses lr$=0.001$, wd$=0.01$, $\beta_{1}=0.9$.  Control best BPB =
$2.6167$; treatment best BPB = $2.7180$; delta = $+0.1013$ BPB
($\sim 3.9\%$ relative, treatment worse).  All data, code, and
ablation design are available on GitHub and archived with Zenodo.
\textbf{Honest caveats:} single seed, single scale, task-specific
benchmark, and no first-principles derivation links $\varphi$ to optimal
learning rates even if the treatment had won.
\end{abstract}

% ═══════════════════════════════════════════════════════════════════════════════
%  INTRODUCTION
% ═══════════════════════════════════════════════════════════════════════════════
\section{Introduction}

The Trinity S$^3$AI program investigates whether geometric invariants of
the H$_4$ Coxeter group (and related structures such as the 600-cell and
Clifford algebra Cl(8)) can encode or constrain the parameters of the
Standard Model of particle physics.  A core epistemic principle of the
program is that \emph{a numeric coincidence is not a derivation}.  This
principle extends to engineering choices: even if a mathematical constant
appears in a hyperparameter recipe, its success or failure on a machine
learning benchmark must be treated as an \texttt{[empirical\_fit]}, not as
a \texttt{[verified]} theorem.

This paper is a \textbf{data paper}.  It documents a controlled ablation
test on the IGLA RACE benchmark; it makes no claim that the result is
derived from first principles, and it reports both the numerical outcome
and its honest limitations.  The canonical claim status for the
$\varphi$-anchored-hyperparameter hypothesis is
\texttt{high\_risk\_or\_falsified} (claim ID:
\texttt{wave24-phi-ablation-falsified}, Layer~L5 of the Trinity claim
ledger).

\paragraph{Prior art.}  The IGLA RACE benchmark (Trinity Waves~18--23)
establishes a character-level language-modeling task on $\varphi$-structured
data and uses Bits Per Byte (BPB) as the primary metric.  Fleet results
(spanning fp8, fp16, bf16, gf16, gf20 precisions and multiple seeds)
report a best-known BPB of $2.4872$ (scarab-fp16-seed77, conventional
hyperparameters).

\paragraph{Gap.}  No controlled ablation had previously tested whether
$\varphi$-anchored hyperparameters outperform conventional values on this
benchmark.

\paragraph{Scope.}  This paper presents exactly such an ablation.  The
negative result closes one speculative avenue (``$\varphi$ hyperparameters
are universally beneficial'') and directs attention back to the physics
content of the project.

% ═══════════════════════════════════════════════════════════════════════════════
%  DATA DESCRIPTION
% ═══════════════════════════════════════════════════════════════════════════════
\section{Data Description}
\label{sec:data}

The ablation compares two arms on an identical model architecture and
training pipeline.  Table~\ref{tab:fixed} lists the fixed parameters.
Table~\ref{tab:arms} lists the hyperparameters that differ between arms.

% ─── Table: Fixed parameters ──────────────────────────────────────────────────
\begin{table}[htbp]
\centering
\caption{Fixed parameters for both ablation arms.}
\label{tab:fixed}
\begin{tabular}{@{}ll@{}}
\toprule
Parameter & Value \\
\midrule
Architecture & 2-layer Transformer (decoder-only) \\
Hidden size & 128 \\
Attention heads & 4 \\
Precision & bf16 \\
Optimizer & AdamW \\
Batch size & 64 \\
Seed & 1597 \\
Total steps & 81\,K \\
Evaluation interval & 10\,K steps \\
Task & Character-level LM on $\varphi$-structured data \\
\bottomrule
\end{tabular}
\end{table}

% ─── Table: Arm hyperparameters ───────────────────────────────────────────────
\begin{table}[htbp]
\centering
\caption{Hyperparameters that differ between the two ablation arms.}
\label{tab:arms}
\begin{tabular}{@{}lcc@{}}
\toprule
Parameter & Control (conventional) & Treatment ($\varphi$-anchored) \\
\midrule
Learning rate & $0.001$ & $\ph^{-3}\approx 0.236$ \\
Weight decay & $0.01$ & $\ph^{-3}\approx 0.236$ \\
$\beta_{1}$ & $0.9$ & $1/\ph\approx 0.618$ \\
$\beta_{2}$ & $0.999$ & $0.999$ \\
\bottomrule
\end{tabular}
\end{table}

% ═══════════════════════════════════════════════════════════════════════════════
%  METHODS
% ═══════════════════════════════════════════════════════════════════════════════
\section{Methods}

The primary metric is Bits Per Byte (BPB), evaluated on a held-out
validation set at 10\,K-step intervals.  Both arms start from identical
initial loss ($\sim 7.0$ BPB).  No early stopping is applied; training
runs to the full 81\,K steps.  Data collection is automated via the
\texttt{trios-trainer-igla} gardener pipeline, which appends BPB
trajectories to a canonical harvest log.

\paragraph{Ablation protocol.}  The two arms are launched
simultaneously on identical hardware.  Wall-clock time is recorded for
budget comparison.  The best BPB achieved by each arm at any step is
tracked as the ``best'' metric; the final BPB at 81\,K steps is tracked
as the ``final'' metric.

% ═══════════════════════════════════════════════════════════════════════════════
%  DATA RECORDS
% ═══════════════════════════════════════════════════════════════════════════════
\section{Data Records}

Table~\ref{tab:trajectory} presents the step-by-step BPB trajectory for
both arms.  Table~\ref{tab:final} presents the final summary metrics.

% ─── Table: Trajectory ────────────────────────────────────────────────────────
\begin{table}[htbp]
\centering
\caption{BPB trajectory at selected checkpoints.}
\label{tab:trajectory}
\begin{tabular}{@{}cccc@{}}
\toprule
Step & Control BPB & Treatment BPB & Delta \\
\midrule
20\,K & $2.6995$ & $3.0922$ & $+0.3927$ \\
30\,K & $2.6378$ & $3.0509$ & $+0.4131$ \\
40\,K & $2.6088$ & $2.9379$ & $+0.3291$ \\
50\,K & $2.5781$ & $2.8210$ & $+0.2429$ \\
60\,K & $\mathbf{2.6167}^{*}$ & $2.8519$ & $+0.2352$ \\
65\,K & $\mathbf{2.6167}^{*}$ & $2.8148$ & $+0.1981$ \\
75\,K & $\mathbf{2.6167}^{*}$ & $2.7483$ & $+0.1316$ \\
81\,K & $\mathbf{2.6167}^{*}$ & $2.7180$ & $+0.1013$ \\
\bottomrule
\end{tabular}
\begin{flushleft}
\small $^{*}$ Control best achieved at 60\,K steps; no further improvement.
\end{flushleft}
\end{table}

% ─── Table: Final metrics ─────────────────────────────────────────────────────
\begin{table}[htbp]
\centering
\caption{Final summary metrics at 81\,K steps.}
\label{tab:final}
\begin{tabular}{@{}lccc@{}}
\toprule
Metric & Control & Treatment & Delta \\
\midrule
Best BPB & $\mathbf{2.6167}$ & $2.7180$ & $+0.1013$ (treatment worse) \\
Final BPB & $2.6596$ & $2.6981$ & $+0.0385$ \\
nCA (higher is better) & $1.287$ & $0.451$ & --- \\
Wall time & $\sim$16\,567 s & $\sim$16\,653 s & comparable \\
\bottomrule
\end{tabular}
\end{table}

% ═══════════════════════════════════════════════════════════════════════════════
%  TECHNICAL VALIDATION
% ═══════════════════════════════════════════════════════════════════════════════
\section{Technical Validation}

\paragraph{Fleet comparison.}  The best BPB achieved across the full
IGLA RACE fleet (as of 2026-05-31) is $2.4872$
(scarab-fp16-seed77, conventional hyperparameters, fp16 precision).
Both ablation arms are substantially above this value (Control $+0.1295$,
Treatment $+0.2308$), confirming that the ablation model (hidden=128) is
under-parameterized relative to fleet leaders.  The \emph{relative}
comparison between Control and Treatment remains valid because both arms
share the identical under-parameterized architecture.

\paragraph{Limitations.}
\begin{enumerate}
  \item \textbf{Single seed.} Only seed 1597 was tested.  Multi-seed
    variance could shift the classification from
    \texttt{high\_risk\_or\_falsified} to \texttt{NULL} or vice versa.
  \item \textbf{Single scale.}  Hidden size 128 is small; $\varphi$
    benefits might emerge at different scales.
  \item \textbf{Task-specific.}  IGLA RACE is a character-level language
    model on $\varphi$-structured data.  Other tasks (vision, protein
    folding, speech) are untested.
  \item \textbf{No theoretical derivation.}  Even if the treatment arm
    had won, the result would be \texttt{[empirical\_fit]}, not
    \texttt{[verified]}, because no first-principles derivation links
    $\varphi$ to optimal learning-rate or weight-decay schedules.
\end{enumerate}

% ═══════════════════════════════════════════════════════════════════════════════
%  USAGE NOTES
% ═══════════════════════════════════════════════════════════════════════════════
\section{Usage Notes}

\paragraph{Reproduction.}  Training logs are stored in the
\texttt{trios-trainer-igla} repository under
\texttt{.trinity/gardener\_harvest.log}.  The canonical ablation design
document is \texttt{docs/audit/PHI\_ABLATION\_DESIGN.md} in the
\texttt{trinity-s3ai} repository.

\paragraph{Extension.}  Additional seeds or scales can be registered in
the ablation design document.  A multi-seed replication (Wave~24.1) is
optional and not prioritized unless the physics track advances.

\paragraph{Validation suite.}  The Trinity repository includes five
automated validators: anti-numerology gate, admitted-honest counter,
claim-generator check, English-only checker, and Markdown link checker.
All must pass before any commit to the main branch.

% ═══════════════════════════════════════════════════════════════════════════════
%  DATA AVAILABILITY
% ═══════════════════════════════════════════════════════════════════════════════
\section{Data Availability}

\begin{itemize}
  \item \textbf{GitHub:} \url{https://github.com/gHashTag/trinity-s3ai}
  \item \textbf{Zenodo:} DOI [PENDING] \texttt{10.5281/zenodo.XXXXXX}
    (reserve after GitHub release tag).
  \item \textbf{License:} MIT (code) / CC-BY-4.0 (data).
  \item \textbf{Canonical ablation design:}
    \texttt{docs/audit/PHI\_ABLATION\_DESIGN.md}
  \item \textbf{Falsification log:}
    \texttt{derivations/falsifiability/wave24\_phi\_ablation.md}
  \item \textbf{Claim ledger (SSOT):} \texttt{docs/claims.yaml}
\end{itemize}

% ═══════════════════════════════════════════════════════════════════════════════
%  ACKNOWLEDGEMENTS
% ═══════════════════════════════════════════════════════════════════════════════
\section*{Acknowledgements}

We thank the anonymous reviewers of Waves~10--19 for forcing the
introduction of the honesty-tag system, and the operators of the IGLA
RACE fleet for maintaining the benchmarking infrastructure.

% ─── Bibliography (inline) ────────────────────────────────────────────────────
\begin{thebibliography}{9}

\bibitem{Koide:1983}
Y. Koide, ``New Mass Relation for Leptons and the Issue of the
  Heaviest Neutrino Mass,''
  \textit{Phys. Lett. B} \textbf{120} (1983) 161--165.

\bibitem{Connes:1996}
A. Connes, ``Gravity Coupled with Matter and the Foundation of
  Noncommutative Geometry,'' \textit{Commun. Math. Phys.}
  \textbf{182} (1996) 155--176.

\bibitem{TrinityWave17}
Trinity S$^3$AI Collaboration,
  ``Boundary Theorem BT-4: Impossibility of H$_4$ Direct Derivation,''
  negative-result companion paper (Wave 17).
  [PENDING] arXiv:XXXX.XXXXX [hep-th] (2026).

\bibitem{TrinityAudit}
Trinity S$^3$AI Collaboration,
  ``Catalog Audit: Sigma Ranking and Epistemic Classification,''
  \texttt{docs/analysis/} (2025).

\end{thebibliography}

\end{document}
